\documentclass{book}

\title{计算线代}
\newcommand{\booksubtitle}{面向计算的线性代数}
\newcommand{\booklicense}{ALL RIGHT PRESERVED}

\author{张津睿}
% Author subtitle could be a university or a geographical location, for example
\newcommand{\authorsubtitle}{新域, 中国}

% Create convenient commands \booktitle and \bookauthor
\makeatletter
\newcommand{\booktitle}{\@title}
\newcommand{\bookauthor}{\@author}
\makeatother

% This utf8 declaration is not needed for versions of latex > 2018 but may
% be helpful for older software. Eventually it may not be worth keeping.
\usepackage[utf8]{inputenc}  
\usepackage{fix-cm}  % this package allows large \fontsize
\usepackage{tikz}    % this is for graphics. e.g. rectangle on title page
\usepackage{amsmath} % Used by equations
\usepackage{amssymb}
\usepackage{CJKutf8}
\usepackage[UTF8]{ctex}
\usepackage[center]{titlesec}%chapter1修改为第1章
\usepackage{hyperref}%生成pdf超链接目录
\hypersetup{hidelinks,
	colorlinks=true,
	allcolors=black,
	pdfstartview=Fit,
	breaklinks=true
}
\titleformat{\chapter}{\raggedright\Huge\bfseries}{第\,\thechapter\,章}{1em}{}
\titleformat{\section}{\raggedright\Large\bfseries}{\,\thesection\,}{1em}{}
\titleformat{\subsection}{\raggedright\large\bfseries}{\,\thesubsection\,}{1em}{}

% The following dimensions specify 4.75" X 7.5" content on 6 3/8" by 9 1/4"
% paper. The paper width and height can be tweaked as required and the content
% should size to fit within the margins accordingly.
%
% The (inside) bindingoffset should be larger for books with more pages. Some
% standard recommended sizes are .375in minimum up to 1in for 600+ page books.
% Sizes .75in and .875in are also recommended roughly at 150 and 400 pages.
\usepackage[bindingoffset=0.625in,
            left=.5in, right=.5in,
            top=.8125in, bottom=.9375in,
            paperwidth=6.375in, paperheight=9.25in]{geometry}
% Here is an alternative geometry for reading on letter size paper:
% \usepackage[margin=.75in, paperwidth=8.5in, paperheight=11in]{geometry}

\renewcommand{\contentsname}{Table of Contents} % default is {Contents}
\usepackage{makeidx}
\makeindex % Initialize an index so we can add entries with \index

% The next few commands are for creating fake content to fill out the template.
% You should delete this (e.g.  everything up to, but not including,
% \begin{document}) after you insert your own content.
% Example content from Einstein's Meaning of Relativity.
% Public domain book: http://www.gutenberg.org/ebooks/36276
\newcommand{\fakeparagraph}{}
\newcommand{\para}{defsdf}
\newcommand{\fakecontent}{
    sldkfjsldkfj

    sdfsdfsdf
}
\newcommand{\prefacePara}
{
现今市面上的线性代数教材非常充足，像我这样一个人
来写一本教材必然会贻笑大方。但是我并不能放下那些
我有时候自己脑海中的古怪想法和天马行空的路线。我
是看面向算子的线性代数入门的，我也感觉非常庆幸能
够在学生时代读到那样的好书。

这本书将基于一种面向计算的方式来讲解一般意义上的
线性代数。线性代数，线性是一个非常广阔的话题，我
非常惶恐地尝试将我在这片非常神奇的领域所学到的东
西分享在这本书中。

这本书会回答所有我曾经感到非常困惑和未能够得到任
意一位老师满意答案的问题。包括行列式为什么这么定
义行向量和列向量到底有什么区别，转置究竟在干什么
，为什么转置有和逆一样的性质（乘积的转置等于分别
转置但是交换相乘），迹又是什么，为什么单单把对角
线元素加起来就得到迹，这种搞笑的计算方法为什么有
那么多神奇的性质，都是那里来的？

In this book, the main goal is to answer all the
questions about those are not well pondered in the
traditional linear algebra course of any kind.
In the computational matrix theory there are plenty
ways of decompose the matrix, such as LU, QR, and the
most fundamental SVD decomposition.
While in the more math oriented linear course
the main topic are always in the most abstract way
as they move forward.
And the main concepts such as the Transpose, dule
space, determinant, trace and all other fundamental
but always treated just in a forced memorize level.

The main goal is to explain all the concepts in a 
more mathematical natrual and fun way, even the sudden
expose of a certain defination would be a very hard
problem to those first expose to the abstract algebra
world student just like me. So I'll try my best to
get every defination and notion in a smooth and more
reasonalbe way. 
}

\newcommand{\fakesections}{
\fakecontent{}
\section{Tensor oriented for computation}\fakecontent{}
\subsection{Examples}\fakecontent{}
\subsection{Excessive Elaborations}\fakecontent{}
\subsection{Long Winded Conclusion}\fakecontent{}
\subsection{Exercises}\fakecontent{}
\section{Theory vs Practice}\fakecontent{}
\subsection{Examples}\fakecontent{}
\subsection{Excessive Elaborations}\fakecontent{}
\subsection{Long Winded Conclusion}\fakecontent{}
\subsection{Exercises}\fakecontent{}
\foreach \x in {A,B,C,D,E,F,G,H,I,J,K,L,M,N,O,P,Q,R,S,T,U,V,W,X,Y,Z}
    {\index{\x 1}\index{\x 2}}
}

% Content Starts Here
\begin{document}
\begin{CJK}{UTF8}{gbsn}
    \frontmatter

    % ---- Half Title Page ----
    % current geometry will be restored after title page
    \newgeometry{top=1.75in,bottom=.5in}
    \begin{titlepage}
        \begin{flushleft}

            % Title
            \textbf{\fontfamily{qcs}\fontsize{48}{54}\selectfont 计算线代\\}

            % Draw a line 4pt high
            \par\noindent\rule{\textwidth}{4pt}\\

            % Subtitle
            % Shaded box from left to right. Text node is midway (centered).
            \begin{tikzpicture}
                \shade[bottom color=lightgray,top color=white]
                (0,0) rectangle (\textwidth, 1.5)
                node[midway] {\textbf{\large {\booksubtitle}}};
            \end{tikzpicture}

            % Edition Number
            \begin{flushright}
                \Large 第一版
            \end{flushright}

            % \vspace{\fill}
            \vspace{\fill}

        \end{flushleft}
    \end{titlepage}
    \restoregeometry
    % ---- End of Half Title Page ----


    % No page numbers on the Frontispiece page
    \thispagestyle{empty}


    % ---- Title Page ----
    % current geometry will be restored after title page
    \newgeometry{top=1.75in,bottom=.5in}
    \begin{titlepage}
        \begin{flushleft}

            % Title
            \textbf{\fontfamily{qcs}\fontsize{48}{54}\selectfont 计线\\}

            % Draw a line 4pt high
            \par\noindent\rule{\textwidth}{4pt}\\

            % Shaded box from left to right with Subtitle
            % The text node is midway (centered).
            \begin{tikzpicture}
                \shade[bottom color=lightgray,top color=white]
                (0,0) rectangle (\textwidth, 1.5)
                node[midway] {\textbf{\large {\booksubtitle}}};
            \end{tikzpicture}

            % Edition Number
            \begin{flushright}
                \Large 第一版
            \end{flushright}

            \vspace{\fill}

            % Author and Location
            \textbf{\large \bookauthor}\\[3.5pt]
            \textbf{\large {\authorsubtitle}}

            \vspace{\fill}

            % Self Publishing Logo. Free to use: CC0 license.
            % The source file is book.svg. If you change the svg, you must then convert
            % it to pdf. There are many online and offline tools available to do that.
            \begin{center}
                \includegraphics{booksvg.pdf}\\[4pt]
                \fontfamily{lmtt}\small{Self Publishers Worldwide\\
                    Seattle San Francisco New York\\
                    London Paris Rome Beijing Barcelona}
            \end{center}

        \end{flushleft}
    \end{titlepage}
    \restoregeometry
    % ---- End of Title Page ----

    % Do not show page numbers on colophon page
    \thispagestyle{empty}

    \begin{flushleft}
        \vspace*{\fill}
        This book was typeset using \LaTeX{} software.\\
        本书给那些爱我与恨我的人。\\
        \vspace{\fill}
        Copyright \textcopyright{} \the\year{}  \bookauthor\\
        License: \booklicense
    \end{flushleft}

    % A title page resets the page # to 1, but the second title page
    % was actually page 3. So add two to page counter.
    \addtocounter{page}{2}

    % The asterisk excludes chapter from the table of contents.
    \chapter{-1 前言}
    \prefacePara{}
    % \[
    %     \left.
    %     \begin{aligned}
    %          & \sqrt{{dX_{1}}^{2} + {dX_{2}}^{2} + {dX_{3}}^{2}} \\
    %          & \qquad
    %         \begin{aligned}
    %                & = \left(1 + \frac{\kappa}{8\pi} \int \frac{\sigma\, dV_{0}}{r}\right)
    %             \sqrt{{dx_{1}}^{2} + {dx_{2}}^{2} + {dx_{3}}^{2}},                             \\
    %             dT & = \left(1 - \frac{\kappa}{8\pi} \int \frac{\sigma\, dV_{0}}{r}\right) dl.
    %         \end{aligned}
    %     \end{aligned}
    %     \right\}
    % \]

    % Three-level Table of Contents
    \setcounter{tocdepth}{3}
    \tableofcontents

    \mainmatter

    \chapter{0 杂谈}

    % 每次阉牛我都在场。对于一般的公牛，只用刀割去即可。但是对于格外生性者，就须采取槌骟术，也就是割开阴囊，掏出睾丸，一木槌砸个稀烂。从此后手术者只知道吃草干活，别的什么都不知道，连杀都不用捆。

    % 后来我才知道生活就是一个缓慢受锤的过程，人一天天老下去，奢望也一天一天消失，之后变得像挨了槌的牛一样。
    % ——黄金时代

    数学并不能防止你受锤，但是在你的黄金时代
    认真地学习数学能防止你睡着，让你醒着受锤，清醒地感受疼痛，和绝望。但同时
    也给予你正确认识和描述世界的语言与改变世界的工具，希望你能用你手里的
    利剑斩开你的前路，希望创造，毁灭的力量就在你的手里了。

    本章并不会证明地讲解若干中学代数之后可能继续发展代数的方向。

    这里面主要还是以坐标计算为主，希望开发若干和高维空间解析几何和物理
    有密切关系的代数运算方法。主要以计算方法为主，并不会牵涉
    公理化的进世代代数方法。

    \section{张量计算方法}

    \subsection{(a,b)-张量 \texorpdfstring{$\forall a,b \in \mathbb{N}$}{∀a,b∈N}的例子}

    $\left[\begin{matrix}
                1 & 2 \\
                3 & 4
            \end{matrix}\right]$
    是一个(1,1)-张量，有一个协变部分和一个逆变部分，
    这种张量其实就是矩阵。矩阵是线性变换的表示。

    $\left[\begin{matrix}
                1 \\
                3
            \end{matrix}\right]$
    是一个(0,1)-张量，这种张量有一个逆变部分，也就是
    之后线性代数中的列向量。

    $\left[\begin{matrix}
                1 & 3
            \end{matrix}\right]$
    是一个(1,0)-张量，这种张量有一个协变部分，也是之后
    线代中的协向量，对偶向量，行向量，线性泛函。

    \subsection{(0,1)-张量 列向量}
    列向量没什么好说的，你写一个大括号，中间放一列东西就行，也许不是数也行。
    $\left[\begin{matrix}
                a \\
                b \\
                c
            \end{matrix}\right]$

    这样的向量就是列向量。

    这样的一列数是没有意义的。聚个例子，我有三个苹果一个
    桃子可以记作
    $\left[\begin{matrix}
                3 \\
                1
            \end{matrix}\right]$
    同时我要是有三个房子一个汽车，我也可以记成
    $\left[\begin{matrix}
                3 \\
                1
            \end{matrix}\right]$
    要是没有单位或者说“基底”来确定每一个分量代表的意义
    这个一列数就没啥意义（或者至少在纯粹的数学以外，很难
    找到意义）

    所以这两个
    $\left[\begin{matrix}
                3 \\
                1
            \end{matrix}\right]$
    应该写成一些基向量的多少倍的和的形式。
    比如把苹果作为基底$\vec{e_1}$，桃子作为基底$\vec{e_2}$,
    房子记作基底$\vec{\tilde{e_1}}$，汽车记作基底$\vec{\tilde{e_2}}$。
    对第一种情况就是$3\vec{e_1}+1\vec{e_2}$
    对第二种情况就是$3\vec{\tilde{e_1}}+1\vec{\tilde{e_2}}$
    我们先把这个东西记成一个更简洁的形式。
    第一种情况
    $\left[\begin{matrix}
                \vec{e_1} & \vec{e_2}
            \end{matrix}\right]
        \left[\begin{matrix}
                3 \\
                1
            \end{matrix}\right]=3\vec{e_1}+1\vec{e_2}$
    第二种情况
    $\left[\begin{matrix}
                \vec{\tilde{e_1}} & \vec{\tilde{e_2}}
            \end{matrix}\right]
        \left[\begin{matrix}
                3 \\
                1
            \end{matrix}\right]=3\vec{\tilde{e_1}}+1\vec{\tilde{e_2}}$
    这两个列向量的分量一样，但是因为基底不同可以有
    不同的意义。


    textbook in linear algebra is very abundant,
    include the linear operator oriented approach,
    which is the currently best one I have ever seen.

    This book will base on a computational oriented
    approach to the subject of the so called linear
    algebra.Linear algebra is such a vast topic that
    I'm having the most firghtened heart to output
    what I have learnt so far in this amazing land.

    In this book, the main goal is to answer all the
    questions about those are not well pondered in the
    traditional linear algebra course of any kind.
    In the computational matrix theory there are plenty
    ways of decompose the matrix, such as LU, QR, and the
    most fundamental SVD decomposition.
    While in the more math oriented linear course
    the main topic are always in the most abstract way
    as they move forward.
    And the main concepts such as the Transpose, dule
    space, determinant, trace and all other fundamental
    but always treated just in a forced memorize level.

    The main goal is to explain all the concepts in a
    more mathematical natrual and fun way, even the sudden
    expose of a certain defination would be a very hard
    problem to those first expose to the abstract algebra
    world student just like me. So I'll try my best to
    get every defination and notion in a smooth and more
    reasonalbe way.
    \subsection{(1,0)-张量 行向量}
    \subsection{(1,1)-张量 线性变换}
    \subsection{(2,0)-张量 双线性函数}
    \subsection{练习}
    \section{外代数计算方法}
    \subsection{几何代数举例}
    \subsection{二向量和三向量}
    \subsection{点积楔积和三维向量场的通量}
    \subsection{练习}
    \foreach \x in {A,B,C,D,E,F,G,H,I,J,K,L,M,N,O,P,Q,R,S,T,U,V,W,X,Y,Z}
        {\index{\x 1}\index{\x 2}}

    \chapter{1 张量(向量) 空间}

    就像「0 杂谈」中说过的那样，单独
    写一个堆有一定结构的分量组成各种
    类型的张量并没有什么意义，除非我
    们在某种基底的意义下讨论个张量。


    \section{张量基础}
    \subsection{基底, 坐标分量, 协变与逆变}
    \subsection{练习}
    \foreach \x in {A,B,C,D,E,F,G,H,I,J,K,L,M,N,O,P,Q,R,S,T,U,V,W,X,Y,Z}
        {\index{\x 1}\index{\x 2}}

    \chapter{2 使用数组表示}

    \chapter{3 度量与度规张量}

    \chapter{Re:1 向量空间}

    \chapter{Re:2 线性映射与作为表示的矩阵}

    \chapter{Re:3 内(外)积空间}

    这一章我们将会回答那些我曾经非常希望知道答案的
    问题，关于行列式，迹，转置，对偶与各种算子（矩
    阵）分解的相关结论和性质。我希望当看完这本书全
    书，和之后这一章后读者能够有一种近乎直观的感受
    感受到这些非常繁琐繁多的性质其实可以非常显然与
    直观。

    \section{内积，作为一种双线性函数}
    \subsection{内积，带度量的空间}

    $\left[\begin{matrix}
                1 & 2 \\
                3 & 4
            \end{matrix}\right]$
    是一个(1,1)-张量，有一个协变部分和一个逆变部分，
    这种张量其实就是矩阵。矩阵是线性变换的表示。

    $\left[\begin{matrix}
                1 \\
                3
            \end{matrix}\right]$
    是一个(0,1)-张量，这种张量有一个逆变部分，也就是
    之后线性代数中的列向量。

    $\left[\begin{matrix}
                1 & 3
            \end{matrix}\right]$
    是一个(1,0)-张量，这种张量有一个协变部分，也是之后
    线代中的协向量，对偶向量，行向量，线性泛函。

    \subsection{内积作为度量}
    列向量没什么好说的，你写一个大括号，中间放一列东西就行，也许不是数也行。
    $\left[\begin{matrix}
                a \\
                b \\
                c
            \end{matrix}\right]$

    这样的向量就是列向量。

    这样的一列数是没有意义的。聚个例子，我有三个苹果一个
    桃子可以记作
    $\left[\begin{matrix}
                3 \\
                1
            \end{matrix}\right]$
    同时我要是有三个房子一个汽车，我也可以记成
    $\left[\begin{matrix}
                3 \\
                1
            \end{matrix}\right]$
    要是没有单位或者说“基底”来确定每一个分量代表的意义
    这个一列数就没啥意义（或者至少在纯粹的数学以外，很难
    找到意义）

    所以这两个
    $\left[\begin{matrix}
                3 \\
                1
            \end{matrix}\right]$
    应该写成一些基向量的多少倍的和的形式。
    比如把苹果作为基底$\vec{e_1}$，桃子作为基底$\vec{e_2}$,
    房子记作基底$\vec{\tilde{e_1}}$，汽车记作基底$\vec{\tilde{e_2}}$。
    对第一种情况就是$3\vec{e_1}+1\vec{e_2}$
    对第二种情况就是$3\vec{\tilde{e_1}}+1\vec{\tilde{e_2}}$
    我们先把这个东西记成一个更简洁的形式。
    第一种情况
    $\left[\begin{matrix}
                \vec{e_1} & \vec{e_2}
            \end{matrix}\right]
        \left[\begin{matrix}
                3 \\
                1
            \end{matrix}\right]=3\vec{e_1}+1\vec{e_2}$
    第二种情况
    $\left[\begin{matrix}
                \vec{\tilde{e_1}} & \vec{\tilde{e_2}}
            \end{matrix}\right]
        \left[\begin{matrix}
                3 \\
                1
            \end{matrix}\right]=3\vec{\tilde{e_1}}+1\vec{\tilde{e_2}}$
    这两个列向量的分量一样，但是因为基底不同可以有
    不同的意义。


    textbook in linear algebra is very abundant,
    include the linear operator oriented approach,
    which is the currently best one I have ever seen.

    This book will base on a computational oriented
    approach to the subject of the so called linear
    algebra.Linear algebra is such a vast topic that
    I'm having the most firghtened heart to output
    what I have learnt so far in this amazing land.

    In this book, the main goal is to answer all the
    questions about those are not well pondered in the
    traditional linear algebra course of any kind.
    In the computational matrix theory there are plenty
    ways of decompose the matrix, such as LU, QR, and the
    most fundamental SVD decomposition.
    While in the more math oriented linear course
    the main topic are always in the most abstract way
    as they move forward.
    And the main concepts such as the Transpose, dule
    space, determinant, trace and all other fundamental
    but always treated just in a forced memorize level.

    The main goal is to explain all the concepts in a
    more mathematical natrual and fun way, even the sudden
    expose of a certain defination would be a very hard
    problem to those first expose to the abstract algebra
    world student just like me. So I'll try my best to
    get every defination and notion in a smooth and more
    reasonalbe way.
    \subsection{练习}
    \section{外积作为度量的可能性}
    \subsection{几何代数举例}
    \subsection{二向量和三向量}
    \subsection{点积楔积和三维向量场的通量}
    \subsection{练习}
    \foreach \x in {A,B,C,D,E,F,G,H,I,J,K,L,M,N,O,P,Q,R,S,T,U,V,W,X,Y,Z}
        {\index{\x 1}\index{\x 2}}

    % \chapter{Determinants}
    % \chapter{Quantics}
    % \chapter{Calculus}
    % \chapter{Differential Equations}
    % \chapter{Infinite Series}
    % \chapter{Theory of Functions}
    % \chapter{Probabilities and Least Squares}
    % \chapter{Analytic Geometry}
    % \chapter{Modern Geometry}
    % \chapter{Trigonometry and Elementary Geometry}
    % \chapter{Non-Euclidean Geometry}
    % \chapter{Bibliography}
    % \chapter{General Tendencies}

    \backmatter
    \addcontentsline{toc}{chapter}{Index}
    \printindex
\end{CJK}
\end{document}
